% =====================================================================
%  AUTHOR'S NOTES — Section 4 of the aluminium tariff report
%  A working companion: where every number in every figure comes from,
%  with the arithmetic shown so it can be checked by hand.
%  Standalone document; not part of the report.
% =====================================================================
\documentclass[11pt, a4paper]{article}
\usepackage[a4paper, top=2.2cm, bottom=2.2cm, left=2.1cm, right=2.1cm]{geometry}
\usepackage{fontspec}
\IfFontExistsTF{DejaVu Sans}
  {\setmainfont{DejaVu Sans}}
  {\IfFontExistsTF{Verdana}{\setmainfont{Verdana}}{\setmainfont{Arial}}}
\usepackage{booktabs}
\usepackage{amsmath}
\usepackage{enumitem}
\usepackage{xcolor}
\usepackage{titlesec}
\usepackage{float}
\usepackage{longtable}
\usepackage{tcolorbox}
\usepackage{hyperref}
\hypersetup{colorlinks=true, linkcolor=black, urlcolor=black}
\emergencystretch=2.5em
\hbadness=10000

\definecolor{shade}{gray}{0.945}
\newtcolorbox{work}[1][Worked example]{colback=shade, colframe=black!50,
  boxrule=0.4pt, arc=2pt, left=6pt, right=6pt, top=4pt, bottom=4pt,
  title={\small\bfseries #1}, coltitle=white, colbacktitle=black!58}
\newtcolorbox{caution}{colback=white, colframe=black!70, boxrule=0.8pt,
  arc=0pt, left=6pt, right=6pt, top=4pt, bottom=4pt}

\titleformat{\section}{\Large\bfseries}{\thesection}{0.7em}{}
\titleformat{\subsection}{\large\bfseries}{\thesubsection}{0.7em}{}
\titleformat{\subsubsection}{\normalsize\bfseries}{\thesubsubsection}{0.7em}{}
\setlength{\parskip}{0.35em}

% \textrm so that these stay upright when used inside math mode
\newcommand{\lakh}{\,\textrm{lakh}}
\newcommand{\cro}{\,\textrm{crore}}

\title{\textbf{Author's Notes: Section 4}\\
\large Every calculation behind every figure and table, with the arithmetic shown}
\author{}
\date{}

\begin{document}
\maketitle
\vspace{-1.0cm}

\begin{center}
\small\itshape A working companion to Section~4 of \textit{Rationalizing Tariff
Structures in the Indian Downstream Aluminium Industry}. It exists so that any
number in the section can be traced to a published ASI cell and re-derived by hand.
Appendix~B of the report states the formal method; this document shows the sums.
\end{center}

\vspace{0.2cm}
\tableofcontents
\newpage

\section{How to reproduce everything}

\subsection{The pipeline}

Three files do all the work.

\begin{table}[H]
\centering
\small
\begin{tabular}{p{4.4cm}p{11.0cm}}
\toprule
\textbf{File} & \textbf{Role} \\
\midrule
\texttt{ASI\_National.csv} & Source data. Long format: one row per
(\texttt{year}, \texttt{nic\_code}, \texttt{indicator}, \texttt{value}). All-India,
sector ``Combined''. Values in INR lakh. \\
\texttt{data\_nic\_ectr.csv} & Effective corporate tax rate by year. One rate per
year, repeated across NIC codes. \\
\texttt{asi\_downstream\_logic.py} & Computes every series, prints the pgfplots
coordinate strings, and writes \texttt{asi\_downstream\_output.txt}. \\
\bottomrule
\end{tabular}
\end{table}

\noindent Running \texttt{python asi\_downstream\_logic.py} regenerates every number
in Section~4. The coordinate strings it prints under heading~8 are pasted directly
into the figures, so a figure cannot silently drift from the data. To rebuild the
report itself: \texttt{python build\_downstream\_report.py} then three passes of
\texttt{xelatex aluminium\_report\_downstream.tex}.

\subsection{Conventions that apply everywhere}

\begin{enumerate}[label=\textbf{C\arabic*.}, leftmargin=2.5em, itemsep=2pt]
\item \textbf{Units.} The ASI publishes in INR lakh. Everything is divided by 100 to
give INR crore for reporting. Employment counts are persons, unconverted.
\item \textbf{Segments.} Upstream $U = \{2420\}$; downstream
$D = \{2432,\ 2511,\ 2599\}$. Aggregation is a plain sum of the class values:
$X_{St} = \sum_{i \in S} X_{it}$. There is no weighting anywhere.
\item \textbf{Years.} Thirteen years are available, 2011--12 to 2023--24. \textbf{All
stated statistics use the twelve years to 2022--23.} 2023--24 appears only in
Figures~4.1, 4.2 and (in earlier drafts) 4.7, marked provisional.
\item \textbf{Ratios, not levels.} Every headline number is a ratio, a share, or a
per-unit figure. This is deliberate and is forced by the apportionment argument
(§\ref{sec:alpha}).
\item \textbf{No deflator.} Everything is at current prices.
\end{enumerate}

\subsection{The variables: what each one is, and what it is not}
\label{sec:vars}

Eight ASI variables carry the whole section. The ASI's definitions are not always what
the plain English name suggests, and three of the eight have traps in them.

\begin{table}[H]
\centering\small
\begin{tabular}{p{1.1cm}p{4.0cm}p{9.6cm}}
\toprule
\textbf{Sym.} & \textbf{ASI indicator} & \textbf{Definition, and what to watch} \\
\midrule
$O_{it}$ & Total Output & Ex-factory value of products and by-products, plus other
receipts. In INR lakh. \\
\addlinespace
$M_{it}$ & Total Inputs & Materials, fuel, electricity and industrial and
non-industrial services consumed. \textbf{Trap:} this is a value, not a quantity, and
it is not broken down by material --- which is why the metal share $\theta$ of the
downstream input bill cannot be observed. \\
\addlinespace
$V_{it}$ & Gross Value Added & Published, and equal to $O_{it} - M_{it}$ to within
rounding. We read the published figure and use the identity only as a check. \\
\addlinespace
$W_{it}$ & Wages and Salaries & Wages to workers plus salaries to other employees.
\textbf{Excludes} employer contributions to provident fund and welfare, which the ASI
reports separately. The labour share $\omega$ is therefore a \emph{lower bound} on
total labour compensation. \\
\addlinespace
$L_{it}$ & Total Number of Persons Engaged & The widest count: all workers, plus
employees other than workers, plus unpaid family members, proprietors and
co-operative members. \textbf{This is the denominator of the headline ratio}, and its
breadth is deliberate --- a narrower count would flatter the downstream case. \\
\addlinespace
$N_{it}$ & Total Number of Workers & Those directly engaged in the manufacturing
process. \textbf{Trap:} $N_{it} < L_{it}$ always. Figure~4.5 uses $N$, not $L$,
because the contract/direct distinction only exists within $N$. \\
\addlinespace
$C_{it}$ & No.\ of Workers employed Through Contractors & Subset of $N_{it}$. The
complement, directly employed workers, is published separately and the two must sum
to $N_{it}$ --- the identity used to reject 2023--24. \\
\addlinespace
$\pi_{it}$ & Net Profit & Gross value added less depreciation, rent, interest and
net indirect taxes. \textbf{Trap:} it is an accounting residual and can be negative,
which is why the tax estimator carries an indicator function. \\
\bottomrule
\end{tabular}
\end{table}

\noindent The ASI also publishes Net Value Added, Net Income, and the split of
employees other than workers into supervisory and other. \textbf{None of these is used
in Section~4.} Net Value Added was considered and rejected: it nets off depreciation,
which is far larger in a capital-intensive smelter than in a fabrication shop, and
using it would have flattered the downstream case by construction. Gross value added
is the conservative choice.

\subsection{The coefficients: five numbers that are not ASI data}

Everything else in Section~4 is a ratio of published cells. Five quantities are not,
and each is a place where the analysis depends on something outside the source. They
are listed here together because they are the section's real exposure.

\begin{table}[H]
\centering\small
\begin{tabular}{p{1.0cm}p{3.4cm}p{10.3cm}}
\toprule
\textbf{Sym.} & \textbf{What it is} & \textbf{Where it comes from, and how much it matters} \\
\midrule
$\tau_t$ & Effective corporate tax rate, by year & From
\texttt{data\_nic\_ectr.csv}; ranges 21.10--28.00 per cent over the period. Applied
uniformly across classes within a year because class-specific rates are not
observable. \textbf{Matters little:} $\tau_t$ cancels exactly from the ratio
$y_D/y_U$, so the employment-yield comparison is unaffected by its level. It affects
only the reported tax \emph{levels}. \\
\addlinespace
$\alpha_i$ & Aluminium share of each NIC class & \textbf{Never estimated, and never
used.} No apportionment is applied. §\ref{sec:alpha} proves the labour-intensity
ratio is bounded within $[3.01, 3.38]$ for \emph{any} positive $\alpha$, so the
quantity does not need to be known. \textbf{Matters not at all} for the ratios;
matters completely for levels, which is why levels are never compared. \\
\addlinespace
$\theta$ & Metal share of the downstream input bill & \textbf{Assumed, not measured.}
The ASI records $M_{it}$ in value only. The report uses 60 per cent illustratively
and tabulates 30--70 per cent. \textbf{Matters a great deal:} it scales the
amplification result linearly, so the ``roughly a fifth to a quarter of value added''
claim is conditional on it. This is the weakest link in the strongest quantitative
claim in the section. \\
\addlinespace
$Q_S$ & Tonnes processed by each segment & \textbf{Not ASI data.} Upstream is
installed primary capacity for FY24 (Ministry of Mines, 4.17 Mt); downstream is an
industry estimate (2.0--2.5 Mt). \textbf{Matters moderately:} it is used only for the
per-tonne restatement in §4.3, which is reported as a range and flagged as supporting
rather than primary evidence. \\
\addlinespace
$m_S$ & Input-cost share & Not a coefficient in the strict sense --- it is computed
as $M_{St}/O_{St}$ from published cells --- but it is listed here because the
amplification factor $m/(1-m)$ is highly convex in it. At $m = 0.83$ the factor is
4.9; at $m = 0.90$ it would be 9.0. Small errors in $m$ propagate. \\
\bottomrule
\end{tabular}
\end{table}

\begin{caution}
\textbf{The honest summary.} Of the five, only $\theta$ is both assumed and
consequential. $\alpha$ is provably irrelevant to what is reported, $\tau$ cancels
from the comparison, $Q$ affects only a supporting statistic, and $m$ is measured.
If Section~4 has one soft spot, it is $\theta$ --- and the fix is two or three
fabricator cost sheets.
\end{caution}

\subsection{The raw cells behind the whole section}

Everything in Section~4 is built from these eight ASI indicators for four NIC classes.
The two years below are the base and terminal years; the same eight indicators are
read for all thirteen years.

\begin{table}[H]
\centering
\footnotesize
\resizebox{\textwidth}{!}{%
\begin{tabular}{llrrrrrrrr}
\toprule
\textbf{Year} & \textbf{NIC} & \textbf{Output} & \textbf{Inputs} & \textbf{GVA} &
\textbf{Wages} & \textbf{Persons} & \textbf{Workers} & \textbf{Contract} & \textbf{Net profit} \\
\midrule
2011--12 & 2420 & 14{,}637{,}024 & 12{,}847{,}817 & 1{,}789{,}207 & 313{,}145 & 122{,}512 & 92{,}527 & 37{,}228 & 711{,}866 \\
2011--12 & 2432 & 736{,}390 & 623{,}139 & 113{,}251 & 37{,}690 & 29{,}192 & 22{,}741 & 9{,}914 & 22{,}691 \\
2011--12 & 2511 & 5{,}628{,}065 & 4{,}544{,}801 & 1{,}083{,}264 & 250{,}427 & 168{,}308 & 132{,}699 & 70{,}411 & 518{,}662 \\
2011--12 & 2599 & 3{,}974{,}627 & 3{,}229{,}008 & 745{,}619 & 274{,}130 & 213{,}703 & 169{,}574 & 74{,}543 & 209{,}884 \\
\midrule
2022--23 & 2420 & 38{,}203{,}861 & 33{,}937{,}702 & 4{,}266{,}159 & 750{,}330 & 170{,}871 & 137{,}611 & 74{,}449 & 1{,}948{,}532 \\
2022--23 & 2432 & 1{,}701{,}255 & 1{,}458{,}413 & 242{,}842 & 95{,}579 & 32{,}677 & 26{,}102 & 12{,}745 & 65{,}464 \\
2022--23 & 2511 & 8{,}759{,}739 & 7{,}514{,}156 & 1{,}245{,}583 & 507{,}301 & 168{,}698 & 133{,}059 & 69{,}647 & 423{,}157 \\
2022--23 & 2599 & 14{,}168{,}873 & 11{,}491{,}965 & 2{,}676{,}908 & 996{,}866 & 322{,}524 & 254{,}030 & 86{,}441 & 1{,}045{,}549 \\
\bottomrule
\end{tabular}}

\vspace{0.15cm}
{\raggedright\footnotesize All monetary values in INR lakh; employment in persons.
Read directly from \texttt{ASI\_National.csv} with no adjustment of any kind.\par}
\end{table}

\noindent Summing the three downstream classes for 2022--23 gives the aggregates used
throughout:

\begin{table}[H]
\centering
\small
\resizebox{\textwidth}{!}{%
\begin{tabular}{lrrrrrrr}
\toprule
& \textbf{Output} & \textbf{Inputs} & \textbf{GVA} & \textbf{Wages} & \textbf{Persons} & \textbf{Workers} & \textbf{Contract} \\
\midrule
Upstream (INR crore / persons) & 382{,}039 & 339{,}377 & 42{,}662 & 7{,}503 & 170{,}871 & 137{,}611 & 74{,}449 \\
Downstream & 246{,}299 & 204{,}645 & 41{,}653 & 15{,}997 & 523{,}899 & 413{,}191 & 168{,}833 \\
\bottomrule
\end{tabular}}
\end{table}

\begin{work}[Check the downstream sum by hand]
GVA: $242{,}842 + 1{,}245{,}583 + 2{,}676{,}908 = 4{,}165{,}333$ lakh $=$ INR
41{,}653\cro. \\
Persons: $32{,}677 + 168{,}698 + 322{,}524 = 523{,}899$. \\
Contract workers: $12{,}745 + 69{,}647 + 86{,}441 = 168{,}833$.
\end{work}

\newpage
\section{Table 4.1 --- the concordance}

\textbf{No calculation.} This table is a classification decision, not a computation,
and it is the one place where judgement rather than arithmetic determines the answer.

The rule is: a NIC class enters the downstream aggregate if and only if its principal
output falls within HS~7604--7616, the thirteen lines on which Section~3 estimates
effective protection. The rule comes from Corden's definition of an
\emph{activity}, and it is applied without exception.

\begin{caution}
\textbf{Where a critic will push, and the answer.}
\begin{itemize}[leftmargin=1.2em, itemsep=2pt]
\item \emph{``NIC 2420 contains semis, which are downstream products.''} True, and
conceded in the text. The ASI classifies an establishment wholly to its principal
product, so semis made inside an integrated smelter are counted upstream. This
inflates upstream and deflates downstream --- it biases the result \emph{against} the
report's own conclusion. No correction is applied.
\item \emph{``Why is 2432 downstream when the material-flow literature calls it a
separate process step?''} Because a tariff does not distinguish a rolling mill from a
window fabricator. Tested directly in Appendix~B.6: reassigning 2432 upstream moves
the ratio from 3.14 to 2.77 and changes nothing else.
\item \emph{``Why exclude 2710 and 2732?''} Their outputs are HS Chapter~85. Including
them is reported as the wide-basket variant; it moves the ratio to 2.61 --- i.e.\
\emph{against} the wide basket --- so the narrow basket is not chosen to flatter.
\item \emph{``2512 is missing.''} It is, and the report says so. HS~7611 and 7613 have
no ASI counterpart, so the downstream aggregate is an undercount.
\end{itemize}
\end{caution}

\section{Figure 4.1 --- gross value added, index 2011--12 $=$ 100}

\subsection{What is plotted}

The gross value added index for both segments, thirteen years including 2023--24.
The persons-engaged index was plotted alongside it in earlier drafts and has been
dropped: the employment comparison is carried more directly by Figure~4.2, and the
index of a headcount that barely moves is not worth a panel. Its values are
retained below because the growth figures quoted in Section~4.2 use them.

\subsection{Formula}
\[
I_{St} \;=\; 100 \times \frac{V_{St}}{V_{S,\,2011\text{--}12}}
\qquad\text{and}\qquad
I^{L}_{St} \;=\; 100 \times \frac{L_{St}}{L_{S,\,2011\text{--}12}}
\]

\begin{work}
\textbf{Upstream GVA index, 2022--23.} $\dfrac{4{,}266{,}159}{1{,}789{,}207} \times 100
= 238.4$. \\[2pt]
\textbf{Downstream GVA index, 2022--23.} $\dfrac{4{,}165{,}333}{1{,}942{,}134} \times 100
= 214.5$. \\[2pt]
\textbf{Upstream persons index, 2022--23.} $\dfrac{170{,}871}{122{,}512} \times 100
= 139.5$. \\[2pt]
\textbf{Downstream persons index, 2022--23.} $\dfrac{523{,}899}{411{,}203} \times 100
= 127.4$.
\end{work}

\subsection{The plotted series}

\begin{table}[H]
\centering\footnotesize
\resizebox{\textwidth}{!}{%
\begin{tabular}{lrrrrrrrrrrrrr}
\toprule
& \textbf{11--12} & \textbf{12--13} & \textbf{13--14} & \textbf{14--15} & \textbf{15--16} & \textbf{16--17} & \textbf{17--18} & \textbf{18--19} & \textbf{19--20} & \textbf{20--21} & \textbf{21--22} & \textbf{22--23} & \textbf{23--24} \\
\midrule
GVA, up & 100 & 174 & 115 & 177 & 169 & 137 & 208 & 165 & 159 & 208 & 354 & 238 & 278 \\
GVA, dn & 100 & 97 & 104 & 111 & 134 & 137 & 136 & 146 & 156 & 140 & 166 & 215 & 238 \\
Persons, up & 100 & 88 & 88 & 88 & 96 & 92 & 104 & 112 & 113 & 124 & 125 & 140 & 161 \\
Persons, dn & 100 & 97 & 106 & 98 & 105 & 111 & 110 & 102 & 111 & 101 & 110 & 127 & 136 \\
\bottomrule
\end{tabular}}
\end{table}

\subsection{Judgement calls}

\begin{itemize}[leftmargin=1.2em, itemsep=2pt]
\item \textbf{Why 2011--12 as base?} It is the first year in the series. Nothing else
recommends it, and \emph{the choice matters}: upstream GVA is volatile enough that
the sign of the growth differential moves with the base year. This is exactly why the
text also reports growth on three-year endpoint averages (§\ref{sec:3yr}), and why
no claim in Section~4 rests on the single-endpoint comparison alone.
\item \textbf{Why plot 2023--24 at all if it is excluded?} Suppressing it invites the
question of what is being hidden. It is plotted and flagged.
\item \textbf{Why only endpoint labels?} The two series cross at 2016--17, where both
index to 137. Labelling every point stacks two identical numbers.
\end{itemize}

\section{Value added per person engaged --- the input to Figure 4.2}

This series is no longer a figure of its own. It was dropped in the brevity pass
because the figure that follows it is its ratio, and two charts for one finding is
one too many. The arithmetic is kept here because every labour-intensity number in
the section is built from it, and because the levels are still quoted in the text.

\subsection{Formula}
\[
v_{St} \;=\; \frac{V_{St}}{L_{St}} \quad \text{(INR lakh per person, since $V$ is in lakh)}
\]

\begin{work}
\textbf{2022--23.} \\
Upstream: $4{,}266{,}159 \div 170{,}871 = 24.967$ lakh per person. \\
Downstream: $4{,}165{,}333 \div 523{,}899 = 7.951$ lakh per person.
\end{work}

\noindent Note that no unit conversion is needed: GVA in lakh divided by persons gives
lakh per person directly. This is a convenient accident of the ASI's units and is worth
remembering when checking.

\subsection{The series (INR lakh per person)}

\begin{table}[H]
\centering\footnotesize
\resizebox{\textwidth}{!}{%
\begin{tabular}{lrrrrrrrrrrrrr}
\toprule
& \textbf{11--12} & \textbf{12--13} & \textbf{13--14} & \textbf{14--15} & \textbf{15--16} & \textbf{16--17} & \textbf{17--18} & \textbf{18--19} & \textbf{19--20} & \textbf{20--21} & \textbf{21--22} & \textbf{22--23} & \textbf{23--24} \\
\midrule
Upstream & 14.6 & 29.1 & 19.1 & 29.4 & 25.7 & 21.7 & 29.2 & 21.5 & 20.5 & 24.5 & 41.5 & 25.0 & 25.3 \\
Downstream & 4.7 & 4.7 & 4.6 & 5.4 & 6.0 & 5.9 & 5.8 & 6.8 & 6.6 & 6.5 & 7.1 & 8.0 & 8.3 \\
\bottomrule
\end{tabular}}
\end{table}

\subsection{What the shape means}

The upstream series tracks the LME aluminium price; the 2021--22 spike to 41.5 is the
metal-price peak, not a productivity event. The downstream series rises almost
monotonically. That asymmetry is the whole reason the section reports a \emph{range}
for the labour-intensity ratio rather than a single year.

\section{Figure 4.2 --- the labour-intensity ratio}
\label{sec:fig43}

This figure carries most of the section's weight, so it is worth being especially
clear about what it is.

\subsection{Formula}
\[
\lambda_t \;=\; \frac{v_{Ut}}{v_{Dt}}
\;=\; \frac{V_{Ut}/L_{Ut}}{V_{Dt}/L_{Dt}}
\]

\noindent Read it as: \emph{a rupee of value added created downstream sustains
$\lambda_t$ times the employment of a rupee created upstream.} The ratio is written
upstream-over-downstream because upstream productivity is the larger number; this
gives a value above one, which is the more natural direction to read.

\begin{work}
\textbf{2022--23.} $\lambda = 24.967 \div 7.951 = \mathbf{3.140}$.
\end{work}

\subsection{The plotted series}

\begin{table}[H]
\centering\footnotesize
\resizebox{\textwidth}{!}{%
\begin{tabular}{lrrrrrrrrrrrr}
\toprule
& \textbf{11--12} & \textbf{12--13} & \textbf{13--14} & \textbf{14--15} & \textbf{15--16} & \textbf{16--17} & \textbf{17--18} & \textbf{18--19} & \textbf{19--20} & \textbf{20--21} & \textbf{21--22} & \textbf{22--23} \\
\midrule
$\lambda_t$ & 3.09 & 6.18 & 4.12 & 5.48 & 4.24 & 3.71 & 5.02 & 3.17 & 3.08 & 3.75 & 5.81 & 3.14 \\
\bottomrule
\end{tabular}}
\end{table}

\noindent Minimum 3.08 (2019--20), maximum 6.18 (2012--13), \textbf{median 3.93}. The
median is the ordinary median of the twelve values, computed on the sorted series;
with twelve observations it is the mean of the sixth and seventh
$(3.75 + 4.12)/2 = 3.935$, printed as 3.93.

\subsection{Two reference lines}

The dashed line is the median. The dotted line at 2.5 is the lower bound of the
Ministry of Mines' own estimate that downstream activity generates two-and-a-half to
three times the employment of primary production. That line is not our estimate; it is
an external benchmark, and its only role is to show that every one of our twelve
observations clears it.

\begin{caution}
\textbf{The most important thing to understand about this figure.} The peaks are
metal-price years. A high LME price raises upstream value added without changing the
number of people a smelter employs, so $\lambda$ rises. It follows that
$\lambda$ is \emph{lowest} in the years least favourable to the report's argument ---
and 2022--23, the year quoted throughout the section, is one of those years. If
challenged on cherry-picking, this is the answer: the quoted year is near the bottom
of the range, not the top.
\end{caution}

\subsection{The per-tonne version}

The text also reports the same finding per tonne of metal. That number is not from the
ASI --- \textbf{tonnage is not an ASI variable} --- and this must be kept straight.

\[
\frac{L_D/Q_D}{L_U/Q_U}
= \underbrace{\frac{L_D/V_D}{L_U/V_U}}_{\lambda = 3.14}
\times \underbrace{\frac{V_D/Q_D}{V_U/Q_U}}_{1.81}
= 5.7
\]

\begin{work}
Upstream: $170{,}871 \div 4{,}170$ thousand tonnes $= 41.0$ persons per thousand tonnes.
$Q_U$ is installed primary capacity, FY24, from the Ministry of Mines. \\[2pt]
Downstream: $523{,}899 \div 2{,}250 = 232.8$ per thousand tonnes at the midpoint;
$210$ at 2.5 Mt and $262$ at 2.0 Mt. $Q_D$ is an industry estimate, hence the range. \\[2pt]
GVA per tonne: upstream $42{,}662\cro \div 4.17$ Mt $=$ INR $102{,}307$ per tonne;
downstream $41{,}653\cro \div 2.25$ Mt $=$ INR $185{,}124$. Ratio $1.81$. \\[2pt]
Check: $3.14 \times 1.81 = 5.68$, against $232.8 \div 41.0 = 5.68$. Consistent.
\end{work}

\begin{caution}
\textbf{Why the report does not use the 35:1 figure that circulates in industry
discussion.} That figure counts the direct workforce of the primary smelters alone
(about 35{,}000--40{,}000, which is 4.6 times smaller than the ASI's count for
NIC~2420) against the direct, indirect \emph{and unorganised} workforce downstream
(8--10 lakh, about 1.7 times larger than the ASI's registered count). Those two
mismatches applied to the consistent ratio give
$5.7 \times 4.6 \times 1.7 \approx 45$, which is where 35--40 comes from. It is not a
different finding; it is the same finding with the denominator widened on one side and
narrowed on the other. Printing it next to our own 3.14 would be indefensible.
\end{caution}

\section{Figure 4.3 --- wage share of GVA}

\subsection{Formula}
\[
\omega_{St} \;=\; \frac{W_{St}}{V_{St}} \times 100
\]

\begin{work}
\textbf{2022--23.} \\
Upstream: $750{,}330 \div 4{,}266{,}159 = 17.59\%$. \\
Downstream: $(95{,}579 + 507{,}301 + 996{,}866) \div 4{,}165{,}333
= 1{,}599{,}746 \div 4{,}165{,}333 = 38.41\%$. \\[2pt]
Ratio: $38.41 \div 17.59 = 2.18$ --- the ``2.2 times as much to household income''
in the text.
\end{work}

\subsection{The plotted series (per cent)}

\begin{table}[H]
\centering\footnotesize
\resizebox{\textwidth}{!}{%
\begin{tabular}{lrrrrrrrrrrrr}
\toprule
& \textbf{11--12} & \textbf{12--13} & \textbf{13--14} & \textbf{14--15} & \textbf{15--16} & \textbf{16--17} & \textbf{17--18} & \textbf{18--19} & \textbf{19--20} & \textbf{20--21} & \textbf{21--22} & \textbf{22--23} \\
\midrule
Upstream & 17.5 & 10.1 & 17.0 & 12.0 & 13.8 & 18.5 & 14.8 & 20.8 & 22.1 & 17.4 & 11.3 & 17.6 \\
Downstream & 28.9 & 32.7 & 35.3 & 33.8 & 33.7 & 37.8 & 41.6 & 39.3 & 39.5 & 42.3 & 39.8 & 38.4 \\
\bottomrule
\end{tabular}}
\end{table}

\subsection{Why this is the most robust number in the section}

$\omega$ is a ratio of two \emph{monetary} variables from the same establishment
return. It therefore does not depend on how persons engaged are counted, which is the
softest variable in the ASI (treatment of contract workers, working proprietors and
unpaid family members all bear on it). If a critic attacks the employment counts,
the wage-share finding survives untouched. Note also that the upstream troughs
(10.1 in 2012--13, 11.3 in 2021--22) are metal-price peaks: a price spike inflates the
denominator, not the numerator.

\section{Figure 4.4 --- GVA by class}

\subsection{Formula}

The same index as Figure~4.1, applied to individual NIC classes rather than segments:
$I_{it} = 100 \times V_{it}/V_{i,\,2011\text{--}12}$.

\begin{work}
\textbf{NIC 2511, 2022--23.} $1{,}245{,}583 \div 1{,}083{,}264 \times 100 = 115.0$. \\
Eleven nominal years, and the class is 15 per cent above where it started.
\end{work}

\subsection{The plotted series}

\begin{table}[H]
\centering\footnotesize
\resizebox{\textwidth}{!}{%
\begin{tabular}{lrrrrrrrrrrrr}
\toprule
& \textbf{11--12} & \textbf{12--13} & \textbf{13--14} & \textbf{14--15} & \textbf{15--16} & \textbf{16--17} & \textbf{17--18} & \textbf{18--19} & \textbf{19--20} & \textbf{20--21} & \textbf{21--22} & \textbf{22--23} \\
\midrule
2420 & 100 & 174 & 115 & 177 & 169 & 137 & 208 & 165 & 159 & 208 & 354 & 238 \\
2599 & 100 & 134 & 145 & 164 & 157 & 181 & 189 & 217 & 209 & 211 & 270 & 359 \\
2432 & 100 & 111 & 107 & 105 & 129 & 168 & 137 & 218 & 179 & 159 & 218 & 214 \\
\textbf{2511} & \textbf{100} & \textbf{70} & \textbf{76} & \textbf{75} & \textbf{119} & \textbf{103} & \textbf{99} & \textbf{90} & \textbf{116} & \textbf{89} & \textbf{89} & \textbf{115} \\
\bottomrule
\end{tabular}}
\end{table}

\noindent \textbf{NIC 2511 is below 100 in seven of the eleven years after the base}
--- 2012--13, 2013--14, 2014--15, 2017--18, 2018--19, 2020--21 and 2021--22. The four
years above are 2015--16 (118.8), 2016--17 (103.2), 2019--20 (116.3) and 2022--23
(115.0). \emph{An earlier draft of the caption said eight; it was miscounted and has
been corrected.}

\subsection{Why this figure matters more than its size suggests}

It is the only place in Section~4 where the ASI evidence and the ERP evidence meet at
the level of an individual product. NIC~2511 makes HS~7610, which carries among the
highest disprotection in the ERP table. It is also the class that has not grown. That
is a coincidence of \emph{class} rather than of aggregate, which is why it is
harder to dismiss than the segment-level comparison --- and why §4.1 is careful to say
it remains association rather than an identified effect.

\section{Table 4.2 --- percentage change by class}

\subsection{Formula}
\[
g_i \;=\; \left( \frac{X_{i,\,2022\text{--}23}}{X_{i,\,2011\text{--}12}} - 1 \right) \times 100
\]

\begin{work}
\textbf{NIC 2420, GVA.} $4{,}266{,}159 \div 1{,}789{,}207 - 1 = +138.4\%$. \\
\textbf{Downstream total, GVA.} $4{,}165{,}333 \div 1{,}942{,}134 - 1 = +114.5\%$. \\
\textbf{NIC 2420, contract workers.} $74{,}449 \div 37{,}228 - 1 = +100.0\%$. \\
\textbf{Downstream, contract workers.} $168{,}833 \div 154{,}868 - 1 = +9.0\%$. \\
\textbf{NIC 2511, corporate tax.} Base $518{,}662 \times 22.01\% = 114{,}158$ lakh;
terminal $423{,}157 \times 22.64\% = 95{,}803$ lakh; change $-16.1\%$.
\end{work}

\noindent Note that the tax change combines a change in profit \emph{and} a change in
the statutory rate. For NIC~2511 profit fell 18.4 per cent while the rate rose
slightly, giving $-16.1$ per cent. This is worth knowing before quoting the figure.

\section{Figure 4.5 --- contract share of workers}

\subsection{Formula}
\[
c_{St} \;=\; \frac{C_{St}}{N_{St}} \times 100
\]
where $C$ is workers employed through contractors and $N$ is total workers.
\textbf{The denominator is total \emph{workers}, not total persons engaged} ---
persons engaged additionally includes employees other than workers and unpaid family
members, who are not part of the contract/direct distinction.

\begin{work}
\textbf{2022--23.} \\
Upstream: $74{,}449 \div 137{,}611 = 54.10\%$. \\
Downstream: $168{,}833 \div 413{,}191 = 40.86\%$, printed 40.9.
\end{work}

\subsection{The plotted series (per cent)}

\begin{table}[H]
\centering\footnotesize
\resizebox{\textwidth}{!}{%
\begin{tabular}{lrrrrrrrrrrrr}
\toprule
& \textbf{11--12} & \textbf{12--13} & \textbf{13--14} & \textbf{14--15} & \textbf{15--16} & \textbf{16--17} & \textbf{17--18} & \textbf{18--19} & \textbf{19--20} & \textbf{20--21} & \textbf{21--22} & \textbf{22--23} \\
\midrule
Upstream & 40.2 & 41.5 & 52.1 & 42.8 & 43.7 & 43.3 & 42.6 & 46.1 & 40.5 & 47.1 & 51.3 & 54.1 \\
Downstream & 47.6 & 48.2 & 46.1 & 41.5 & 44.1 & 40.5 & 41.5 & 42.0 & 39.3 & 37.1 & 40.6 & 40.9 \\
\bottomrule
\end{tabular}}
\end{table}

\subsection{Why 2023--24 is excluded from this figure specifically}

The ASI publishes directly-employed workers and contract workers separately, and they
should sum to total workers. They do, exactly, in every year from 2011--12 to
2022--23 --- residuals of one or two workers, pure rounding. In 2023--24 the residual
across the eight classes originally examined is \textbf{178{,}591 workers, 19.6 per
cent of the total}. Something changed in the definition or the coverage. The contract
share cannot be computed on that basis and the year is dropped.

\begin{caution}
\textbf{A claim that was removed.} An earlier draft said tariff shocks transmit into
contract-worker adjustment ``with very little lag''. Twelve annual observations
against three years of ERP estimates identify no lag structure, and the series in fact
moves gradually. The claim was removed rather than defended. What replaces it needs no
timing assumption: downstream contract employment grew 9.0 per cent while upstream
doubled --- a composition result, not a dynamic one.
\end{caution}

\section{Figure 4.6 --- employment yield per rupee of estimated tax}

\subsection{This is an estimate, and the distinction matters}

Only the employment yield is plotted. The corporate-tax levels that fed the bar panel
in earlier drafts were dropped: on levels upstream is ahead, so the panel conceded a
point without advancing one, and the fact is now stated in a single sentence of
text instead. The tax estimator is still needed, because the yield is built on it.

The ASI does not publish tax paid. It publishes \emph{net profit}. The tax series is
constructed:
\[
T_{St} \;=\; \tau_t \sum_{i \in S} \pi_{it}\,\mathbf{1}\{\pi_{it} > 0\}
\]
$\tau_t$ is the effective corporate tax rate for the year, from
\texttt{data\_nic\_ectr.csv}. The indicator matters: a loss-making class would
otherwise generate a negative tax, which has no behavioural meaning. One observation
within the retained classes is affected --- NIC~2432 in 2013--14.

\begin{table}[H]
\centering\footnotesize
\resizebox{\textwidth}{!}{%
\begin{tabular}{lrrrrrrrrrrrrr}
\toprule
& \textbf{11--12} & \textbf{12--13} & \textbf{13--14} & \textbf{14--15} & \textbf{15--16} & \textbf{16--17} & \textbf{17--18} & \textbf{18--19} & \textbf{19--20} & \textbf{20--21} & \textbf{21--22} & \textbf{22--23} & \textbf{23--24} \\
\midrule
$\tau_t$ (\%) & 22.01 & 21.10 & 23.20 & 22.06 & 25.86 & 24.75 & 28.00 & 27.21 & 21.95 & 25.68 & 22.64 & 22.64 & 22.32 \\
\bottomrule
\end{tabular}}
\end{table}

\begin{work}[The 2022--23 tax build-up, class by class]
$\tau = 22.64\%$ for all classes. \\[2pt]
2420: $1{,}948{,}532 \times 0.2264 = 441{,}148$ lakh $=$ INR $4{,}411\cro$ \\
2432: $\ \ \ \ 65{,}464 \times 0.2264 = \ \ 14{,}821$ lakh $=$ INR $\ \ 148\cro$ \\
2511: $\ \ 423{,}157 \times 0.2264 = \ \ 95{,}803$ lakh $=$ INR $\ \ 958\cro$ \\
2599: $1{,}045{,}549 \times 0.2264 = 236{,}712$ lakh $=$ INR $2{,}367\cro$ \\[2pt]
Downstream total: $148 + 958 + 2{,}367 = $ INR $3{,}473\cro$.
\end{work}

\subsection{Employment yield --- what is plotted}
\[
y_{St} \;=\; \frac{L_{St}}{T_{St}} \quad \text{(persons per INR crore of estimated tax)}
\]

\begin{work}
\textbf{2022--23.} Upstream $170{,}871 \div 4{,}411 = 38.7$, printed 39. \\
Downstream $523{,}899 \div 3{,}473 = 150.8$, printed 151. Ratio $3.89$.
\end{work}

\begin{table}[H]
\centering\footnotesize
\resizebox{\textwidth}{!}{%
\begin{tabular}{lrrrrrrrrrrrr}
\toprule
& \textbf{11--12} & \textbf{12--13} & \textbf{13--14} & \textbf{14--15} & \textbf{15--16} & \textbf{16--17} & \textbf{17--18} & \textbf{18--19} & \textbf{19--20} & \textbf{20--21} & \textbf{21--22} & \textbf{22--23} \\
\midrule
Tax, up (INR '000 cr) & 1.57 & 4.64 & 2.38 & 4.10 & 3.94 & 2.22 & 5.74 & 2.78 & 1.99 & 4.13 & 9.00 & 4.41 \\
Tax, dn & 1.65 & 1.42 & 1.42 & 1.53 & 2.34 & 1.91 & 1.78 & 2.01 & 1.51 & 1.88 & 2.44 & 3.47 \\
Yield, up & 78 & 23 & 45 & 26 & 30 & 51 & 22 & 49 & 70 & 37 & 17 & 39 \\
Yield, dn & 249 & 281 & 306 & 263 & 184 & 238 & 254 & 208 & 301 & 221 & 185 & 151 \\
\bottomrule
\end{tabular}}
\end{table}

\begin{caution}
\textbf{Limitations to state if asked.} $\tau_t$ is applied uniformly across classes
within a year, because class-specific effective rates are not observable. The
estimate takes no account of exemptions, carried-forward losses, depreciation timing
or minimum alternate tax. It is a proxy for the \emph{tax base} attributable to each
segment, not a record of receipts. The comparison between segments is more reliable
than either level, because both are constructed identically and the same $\tau_t$
cancels from the ratio $y_D/y_U$ entirely.
\end{caution}

\subsection{The volatility statistic quoted alongside}

\[
\sigma_S \;=\; \operatorname{sd}\!\big( \ln V_{S,t+1} - \ln V_{St} \big),
\quad \text{sample standard deviation, } n-1 \text{ denominator, eleven differences}
\]

\begin{table}[H]
\centering\footnotesize
\resizebox{\textwidth}{!}{%
\begin{tabular}{lrrrrrrrrrrr}
\toprule
$\Delta \ln V$ & \textbf{12--13} & \textbf{13--14} & \textbf{14--15} & \textbf{15--16} & \textbf{16--17} & \textbf{17--18} & \textbf{18--19} & \textbf{19--20} & \textbf{20--21} & \textbf{21--22} & \textbf{22--23} \\
\midrule
Upstream & $+0.556$ & $-0.417$ & $+0.432$ & $-0.047$ & $-0.213$ & $+0.420$ & $-0.231$ & $-0.039$ & $+0.269$ & $+0.533$ & $-0.395$ \\
Downstream & $-0.033$ & $+0.072$ & $+0.065$ & $+0.189$ & $+0.021$ & $-0.009$ & $+0.073$ & $+0.064$ & $-0.107$ & $+0.171$ & $+0.256$ \\
\bottomrule
\end{tabular}}
\end{table}

\noindent $\sigma_U = 0.3732$, $\sigma_D = 0.1046$, ratio $\mathbf{3.57}$ (reported as
``3.6 times as volatile''). Log differences are used rather than percentage changes
because they are symmetric in direction and additive over time, which matters when a
series moves by $+74$ and $-34$ per cent in adjacent years.

\section{Figure 4.7 --- the four multiples}

Each bar is a ratio of two numbers already computed elsewhere in the section. Nothing
new is calculated; the figure exists to put four measures drawn from four different
parts of the ASI schedule on one comparable scale.

\begin{table}[H]
\centering\small
\begin{tabular}{llr}
\toprule
\textbf{Measure} & \textbf{Arithmetic} & \textbf{Multiple} \\
\midrule
Labour intensity & $24.967 \div 7.951$ & $3.14$ \\
Labour share of value added & $38.41 \div 17.59$ & $2.18$ \\
Jobs per rupee of estimated tax & $150.83 \div 38.73$ & $3.89$ \\
Volatility & $0.3732 \div 0.1046$ & $3.57$ \\
\bottomrule
\end{tabular}
\end{table}

\noindent The first three are oriented so that a higher number favours downstream on
employment or household-income grounds. \textbf{The fourth is oriented differently and
this must be said when presenting it}: it is upstream volatility divided by downstream,
so a high number means the \emph{upstream} tax base is the less dependable one. The
figure's caption says so; a slide deck reproducing the chart should too.

\section{The three-year endpoint growth rates}
\label{sec:3yr}

Because upstream GVA is volatile, single-endpoint growth is unstable. The text
therefore also reports growth between three-year averages, with
$A = \{2011\text{--}12,\,2012\text{--}13,\,2013\text{--}14\}$ and
$B = \{2020\text{--}21,\,2021\text{--}22,\,2022\text{--}23\}$.

\begin{work}
\textbf{Upstream.} $A = (17{,}892,\ 31{,}191,\ 20{,}552)$, mean $23{,}211.6\cro$;
$B = (37{,}157,\ 63{,}340,\ 42{,}662)$, mean $47{,}719.6\cro$. \\
Growth $= 47{,}719.6/23{,}211.6 - 1 = +105.6\%$;
CAGR $= (2.056)^{1/9} - 1 = 8.34\%$. \\[3pt]
\textbf{Downstream.} $A$ mean $19{,}472.2$; $B$ mean $33{,}687.5$.
Growth $+73.0\%$; CAGR $6.28\%$. \\[3pt]
\textbf{NIC 2511.} $A$ mean $8{,}859.3$; $B$ mean $10{,}574.7$.
Growth $+19.4\%$; CAGR $1.99\%$.
\end{work}

\noindent \textbf{Why the exponent is 9.} The midpoint of $A$ is 2012--13 and the
midpoint of $B$ is 2021--22; nine years separate them. Using 11 (the endpoint span)
would understate the compound rate.

\section{The apportionment argument}
\label{sec:alpha}

This is the most important robustness result in the section and the one most likely to
be challenged, so the logic is worth restating in full.

\textbf{The objection.} None of the four classes is exclusive to aluminium. NIC~2599
in particular contains a great deal of non-aluminium fabrication. Should not each
class be scaled by its aluminium share $\alpha_i$ before aggregating?

\textbf{The answer.} It makes no material difference, and this can be \emph{proved}
rather than simulated. Writing $w_i(\alpha) = \alpha_i L_i / \sum_j \alpha_j L_j$,
\[
v_D(\alpha) = \frac{\sum_i \alpha_i V_i}{\sum_i \alpha_i L_i}
= \sum_i w_i(\alpha)\, \frac{V_i}{L_i} = \sum_i w_i(\alpha)\, v_i ,
\]
a weighted average of the class productivities with positive weights summing to one.
It is therefore bounded by $\min_i v_i$ and $\max_i v_i$ \emph{whatever} $\alpha$ is.
Upstream is a single class, so $\alpha_U$ cancels exactly.

\begin{work}
2022--23 class productivities: $v_{2432} = 7.4316$, $v_{2511} = 7.3835$,
$v_{2599} = 8.2999$ lakh per person; $v_U = 24.9671$. \\[2pt]
Lower bound: $24.9671 \div 8.2999 = 3.008$. \\
Upper bound: $24.9671 \div 7.3835 = 3.382$. \\[2pt]
So $\lambda(\alpha) \in [3.01,\ 3.38]$ for any positive apportionment vector, against
the undiscounted $3.14$. Verified against 200{,}000 random draws of $\alpha$ from
$U(0.01, 1)^3$: observed range $[3.011,\ 3.373]$.
\end{work}

\begin{caution}
\textbf{What is \emph{not} protected.} Levels are homogeneous of degree one in
$\alpha$ and move freely: the cross-segment employment multiple ranges from 0.81 to
3.07 across plausible vectors. This is why Section~4 never compares employment levels
across segments, and why every headline is a ratio. It is a discipline imposed by the
sensitivity analysis, not a stylistic preference.
\end{caution}

\section{The input-cost amplification factor}
\label{sec:amp}

Used in §4.7 to give the transmission mechanism a size.

\[
V = (1-m)O \;\Longrightarrow\;
\frac{\Delta V}{V} = -\,\frac{m\theta}{1-m}\,\hat p
\]
where $m$ is bought-in inputs as a share of output value, $\theta$ the metal share of
the input bill and $\hat p$ the proportional change in the metal price.

\begin{work}
$m_D = 204{,}645 \div 246{,}299 = 0.8309$. Amplification $= 0.8309/0.1691 = 4.91$. \\
At $\theta = 0.6$: a 1\% rise in metal price costs $4.91 \times 0.6 = 2.95\%$ of
downstream value added; the 7.5\% duty costs $7.5 \times 2.95 = 22.1\%$.
\end{work}

\begin{caution}
\textbf{$\theta$ is the one parameter not measured.} The ASI records the input bill in
value but not by material, so the metal share is unobservable. The report uses 60 per
cent illustratively and Appendix~B tabulates 30--70 per cent. \textbf{This is the
weakest link in the strongest quantitative claim in Section~4}, and it is why the
decision record lists fabricator cost sheets as a data priority. Do not quote the
22 per cent figure without the conditional.
\end{caution}

\section{Cross-checks that were run}

\begin{enumerate}[leftmargin=1.6em, itemsep=3pt]
\item \textbf{Accounting identity.} Output $-$ inputs $-$ GVA $= 0$ to within
$10^{-6}$ per cent in every class and every year. Confirms the three variables are
read from the correct rows.
\item \textbf{Worker identity.} Directly employed $+$ contract $=$ total workers,
exactly, in all twelve years to 2022--23 (residual of $\pm 2$ workers, rounding).
Fails in 2023--24 with a residual of 178{,}591.
\item \textbf{Reproduction of the earlier memo.} The pipeline independently
reproduces the figures in \texttt{NIC\_Selection\_Review\_Memo.md} --- 3.14, 38.4 per
cent, 151 against 39, $+138.4$ against $+114.5$ --- which were computed separately.
\item \textbf{Analytic bound against simulation.} The apportionment bound
$[3.01, 3.38]$ was checked against 200{,}000 random draws.
\item \textbf{Per-tonne identity.} $3.14 \times 1.81 = 5.68$ reconciles with
$232.8 / 41.0 = 5.68$ computed directly.
\end{enumerate}

\section{The arguments in favour of the downstream industry, and the figure that carries each}

Section~4 exists to establish one proposition:

\begin{center}
\itshape that the segment the tariff structure disprotects is the segment on which
employment and household income depend.
\end{center}

\noindent Seven figures survive the final draft. Each carries a distinct argument, and
none is a restatement of another --- they draw on value added, employment, wages,
corporate tax and the variance of value added, which are five different parts of the
ASI schedule. This section sets out what each argument claims, what it rests on, how
strong it is, and where it can be attacked. It is written as though by an opponent,
because that is the only useful way to write it.

\subsection{Argument 1 --- A rupee spent downstream buys more employment}
\textbf{Figure 4.2, the labour-intensity ratio. The central argument.}

\textbf{The claim.} Value added created in downstream fabrication sustains
\textbf{3.14 times} the employment of the same value added created in primary metal;
across the record the multiple runs from 3.08 to 6.18 with a median of 3.93.

\textbf{Why it is the central argument.} It is the one that converts a tariff question
into an employment question. A tariff decision reallocates margin between segments;
this ratio says what each rupee of reallocated margin is worth in jobs. Nothing else
in the report does that.

\textbf{What makes it hard to attack.} Twelve years, twelve confirmations, no
exception. A critic must argue not that the number is wrong but that it is wrong
\emph{every year for twelve years in the same direction}. The apportionment objection
--- that the classes are not purely aluminium --- is closed analytically rather than by
assertion: the ratio is bounded within $[3.01, 3.38]$ for \emph{any} positive
apportionment vector (§\ref{sec:alpha}).

\textbf{Where it can be attacked.} The denominator is persons engaged, the softest
variable in the ASI. Anyone disputing the treatment of contract workers, working
proprietors or unpaid family members is disputing this figure. The answer is
Argument~2, which does not use employment counts at all.

\textbf{Presentational note.} The peaks are metal-price years, so the ratio is
\emph{lowest} in the years least favourable to us --- and 2022--23, the year quoted
throughout, is one of them. If accused of cherry-picking, say this.

\subsection{Argument 2 --- More of what downstream earns reaches households}
\textbf{Figure 4.3, wage share of value added. The most robust argument.}

\textbf{The claim.} Downstream converts \textbf{38.4 per cent} of gross value added
into wages and salaries; upstream converts \textbf{17.6 per cent}. A rupee retained
downstream therefore delivers roughly \textbf{2.2 times} as much to household income.

\textbf{Why it matters separately from Argument 1.} Argument 1 counts heads;
this one follows the money. A segment could employ many people badly. This figure
shows that downstream does not merely engage more people per rupee but passes more of
each rupee to them, which is the outcome a policymaker actually cares about.

\textbf{What makes it the most robust.} Both numerator and denominator are monetary
figures from the same establishment return. The figure is therefore immune to every
objection about how employment is counted --- the single largest vulnerability in the
rest of the section. \textbf{If the section had to be reduced to one chart, this is
the one to keep.}

\textbf{Where it can be attacked.} Wages exclude employer contributions to provident
fund and welfare, which the ASI reports separately; the true labour share is higher in
both segments. This understates the level but there is no reason to think it changes
the gap. The two series never cross in twelve years.

\subsection{Argument 3 --- The protected segment is not the one in difficulty}
\textbf{Figure 4.1, gross value added index.}

\textbf{The claim.} Over the eleven years to 2022--23 upstream value added rose
\textbf{138.4 per cent} against \textbf{114.5 per cent} downstream, and upstream
overtook a downstream segment that had started the period larger.

\textbf{What it is for.} It answers the prior question a reader will ask before any of
the rest matters: \emph{if the duty protects upstream, is upstream the segment that
needs protecting?} On this evidence, no. Protection is flowing to the faster-growing
of the two. The figure also establishes that the segments are of comparable size,
without which every subsequent comparison reads as a large thing set against a small
one.

\textbf{Where it can be attacked, and this one is genuinely exposed.} Single-endpoint
growth is base-year sensitive, and upstream value added is volatile enough that the
sign of the differential can be moved by the base chosen. This is why the text also
gives the three-year-average comparison (105.6 against 73.0 per cent), which is stable
and says the same thing. \emph{Never quote the endpoint figure without the three-year
one.}

\textbf{The risk in the figure itself.} Read alone, it shows downstream value added up
114 per cent, which a hostile reader can quote as evidence that the segment is
thriving. The framing sentence must always travel with the chart.

\subsection{Argument 4 --- The damage appears exactly where the tariff analysis predicts}
\textbf{Figure 4.4, value added by class. The most important figure for the thesis.}

\textbf{The claim.} Structural metal products (NIC~2511), the class that makes
HS~7610 --- aluminium frameworks, doors, windows and fa\c{c}ades --- grew nominal value
added by \textbf{15.0 per cent} in eleven years while its workforce did not grow at
all and its estimated tax contribution fell 16.1 per cent. HS~7610 carries among the
highest disprotection in the entire ERP table.

\textbf{Why this is the most important figure.} It is the only place in the report
where Section~3 and Section~4 meet on the \emph{same object}. Everywhere else the
tariff evidence and the industrial evidence run in parallel and the reader is asked to
connect them. Here the connection is made at the level of a single product line: the
line the tariff analysis says is worst treated is the line the industrial statistics
say has not grown. A coincidence of \emph{class} is much harder to dismiss than a
coincidence of aggregates.

\textbf{Where it can be attacked, and the attack is good.} The construction cycle. Structural
products track building activity closely, and the period covers a real-estate downturn
that has nothing to do with tariffs. We do not control for it, and Section~4.1 now
concedes the point explicitly rather than waiting to be caught. The honest position is
that this figure is the strongest \emph{suggestive} evidence in the report and not
proof of anything.

\subsection{Argument 5 --- The harm shows up as jobs never created}
\textbf{Figure 4.5, contract share of workers.}

\textbf{The claim.} Upstream the contract workforce doubled and its share rose from
40.2 to \textbf{54.1 per cent}. Downstream contract employment grew 9.0 per cent and
the share \emph{fell} to \textbf{40.9 per cent}, while downstream value added more
than doubled.

\textbf{The argument, which is subtler than it looks.} A falling contract share
normally reads as improving job quality. Here it means the opposite: downstream value
added doubled while its flexible workforce barely moved, which is the signature of a
segment running existing capacity harder rather than adding to it. The harm is not
redundancy but \emph{foregone hiring} --- and foregone hiring appears in no statistic
that anyone collects.

\textbf{Where it can be attacked.} This is the weakest of the seven, and it got weaker
when the claim that tariff shocks transmit into contract adjustment ``with very little
lag'' was removed --- twelve annual observations against three years of ERP estimates
identify no lag structure. What remains is a description, not a mechanism. Retain it,
but do not lead with it.

\subsection{Argument 6 --- The revenue at risk is thin, and cyclical}
\textbf{Figure 4.6, employment yield per rupee of estimated tax.}

\textbf{The claim.} Downstream supports \textbf{151 persons engaged per crore} of
estimated corporate tax; upstream supports \textbf{39}. The ordering holds in all
twelve years and the gap never falls below a factor of two.

\textbf{What it is for.} It pre-empts the only serious fiscal objection to the
proposal --- that cutting the duty forgoes revenue. The answer is that revenue should be
weighed by what depends on it. Rationalising the duty trades a revenue stream thinly
attached to employment for relief to one on which roughly four times as many
livelihoods depend per rupee at stake. The volatility statistic quoted alongside makes
the same point from the exchequer's side: upstream value added is \textbf{3.6 times}
as variable, so the revenue being defended is also the less dependable of the two.

\textbf{An honest concession that stays in the text.} On \emph{levels}, upstream pays
more --- an estimated INR 4,411 crore against INR 3,473 crore. This was a bar panel in
earlier drafts and was cut, because on levels the comparison favours upstream and the
panel conceded a point without advancing one. The fact remains in a sentence of text.
\textbf{Do not remove that sentence.} A submission that suppresses the one number
running against it invites the reader to distrust the rest.

\textbf{Where it can be attacked.} The tax series is an estimate, not receipts, and
applies a uniform effective rate across classes. Both segments are constructed
identically and the rate cancels from the ratio entirely, so the comparison is far
sounder than either level.

\subsection{Argument 7 --- All of it, on one scale}
\textbf{Figure 4.7, the four multiples.}

\textbf{What it is for.} A recapitulation, and the natural figure for a one-page
brief or a slide. Its value is that the four measures come from four different parts
of the ASI schedule, so it demonstrates that the case does not rest on one statistic
viewed four ways.

\textbf{The one thing to watch.} The fourth bar is oriented differently from the
first three: it is upstream volatility divided by downstream, so a high value means
the \emph{upstream} tax base is the less dependable one. Explain this wherever the
chart appears.

\subsection{The argument with no figure}

The \textbf{transmission mechanism} --- the reason a modest input tariff is not a modest
problem --- has no chart. Bought-in inputs are 83.1 per cent of downstream output value,
so value added is 16.9 per cent, and the ratio of the two, \textbf{4.91}, is the factor
by which any proportional change in input cost is amplified by the time it reaches
value added. Where metal is 60 per cent of the input bill, the 7.5 per cent duty
accounts for something over a fifth of downstream value added; upstream the figure is
nil, because a smelter does not buy the metal it makes.

This is the analytical heart of the report --- it is what makes the tariff argument and
the employment argument the same argument --- and it is carried entirely by two
sentences of text and an appendix derivation. That is a deliberate choice, because the
one parameter it depends on ($\theta$, the metal share of the input bill) is assumed
rather than measured, and a chart would lend it a precision it has not earned. If
fabricator cost sheets are ever obtained, this becomes the strongest quantitative
claim in the report and would then deserve a figure.

\subsection{Ranking, for when there is not room for everything}

\begin{table}[H]
\centering\small
\begin{tabular}{clp{7.4cm}}
\toprule
& \textbf{Argument} & \textbf{Keep because} \\
\midrule
1 & Wage share (Fig.~4.3) & Immune to every objection about employment counts. \\
2 & Labour intensity (Fig.~4.2) & The central finding; twelve-for-twelve. \\
3 & Class disaggregation (Fig.~4.4) & The only link between the tariff and industrial evidence. \\
4 & Employment yield (Fig.~4.6) & Answers the fiscal objection. \\
5 & GVA index (Fig.~4.1) & Establishes comparable size and who is actually growing. \\
6 & Summary multiples (Fig.~4.7) & Presentation, not evidence. \\
7 & Contract share (Fig.~4.5) & Descriptive; the weakest, and dispensable under pressure. \\
\bottomrule
\end{tabular}
\end{table}

\subsection{What is missing}

Three things would strengthen the case and are not in the section.

\begin{enumerate}[leftmargin=1.6em, itemsep=3pt]
\item \textbf{A direct link between effective protection and downstream outcomes.}
The natural figure is ERP against class growth. It cannot be drawn: three years of ERP
estimates against twelve of ASI data would produce a picture that overstates what we
know. This absence is the honest reason Section~4.1 says the evidence is association.
\item \textbf{Input-cost share over time.} It is computed --- downstream rises from 81.2
to 83.1 per cent over the period --- but quoted for a single year only. A figure showing
the squeeze tightening would support the transmission mechanism, which currently has
no chart at all.
\item \textbf{Capacity utilisation.} The most vivid downstream statistic in the report
(30--35 per cent against upstream's near-98 per cent) is not an ASI variable and cannot
be brought into Section~4 on the same footing.
\end{enumerate}

\section{The verification harness}

\texttt{verify\_figures.py} recomputes every series from
\texttt{ASI\_National.csv}, parses every \texttt{\textbackslash addplot} coordinate
block out of the section source, and asks of each plotted number whether it is a
correct rounding of the computed value \emph{at the precision it is printed}. It also
re-checks the thirty headline and derived figures quoted in the prose. Run it after
any edit to the section:

\begin{center}\texttt{python verify\_figures.py}\end{center}

\noindent It currently reports \textbf{ALL CHECKS PASS} across \textbf{158 plotted points} and
30 quoted statistics. It caught four real errors on its first run --- three plotted
points that had been rounded twice (136.47 printed as 137, 214.47 as 215, 139.47 as
140) and the amplification factor, which had been computed from the rounded input
share 0.831 rather than the actual 0.8309 and was printed as 4.92 instead of 4.91.
All four are now fixed. None changed an argument, but the fact that they survived
several drafts is the case for running the harness rather than trusting the prose.

\textbf{Keep the harness in step with the figures.} It matches each computed series
against the coordinate block that plots it, so a series that is no longer plotted will
be matched against the wrong block and reported as a spurious mismatch. When a figure
is added or dropped, update the \texttt{SERIES} dictionary at the top of the script
before reading its output. It was last updated when the section went from eight
figures to seven.

\section{Known weak points}

Stated here so that nobody has to discover them from a hostile reviewer.

\begin{enumerate}[leftmargin=1.6em, itemsep=3pt]
\item \textbf{$\theta$ is assumed, not measured} (§\ref{sec:amp}). The 22 per cent figure is
conditional and must always be presented as such.
\item \textbf{Tonnage is not ASI data.} The per-tonne figures mix official employment
with industry-estimated tonnage. They are supporting evidence, not headline.
\item \textbf{The tax series is a proxy}, with a uniform rate across classes.
\item \textbf{Causation is not established.} The report now says so explicitly in
§4.1. The construction cycle is an uncontrolled alternative explanation for the
NIC~2511 result specifically, and it is a good one --- structural metal products
follow building activity closely.
\item \textbf{The ASI frame excludes unorganised units}, which are concentrated
downstream. This understates downstream employment, so it runs in the report's
favour --- but it also means the downstream aggregate is not the whole sector.
\end{enumerate}

\vspace{0.3cm}
\noindent One error was found and fixed while preparing these notes: the Figure~4.5
caption said NIC~2511 spent \emph{eight} of eleven years below its base level. The
correct count is seven. The figure itself was always right; only the caption was
wrong. It is worth recording that the error survived several drafts and was caught
only by writing the arithmetic out --- which is the argument for this document
existing.

\end{document}
